\section{Implémentation}
\label{sec:impl}

Cette section présente les éléments d'implémentation les plus structurants
du projet. Le code complet est disponible dans le dépôt~; nous nous
limitons ici aux extraits qui répondent directement aux exigences
académiques.

\subsection{Configuration de l'internationalisation}

La gestion des trois langues repose entièrement sur l'infrastructure
\texttt{gettext} de Django.

\begin{lstlisting}[language=Python,caption={Extrait de
crm\_project/settings.py},label=lst:i18n]
from pathlib import Path
from django.utils.translation import gettext_lazy as _

LANGUAGE_CODE = 'en'
LANGUAGES = [
    ('en', _('English')),
    ('fr', _('French')),
    ('ar', _('Arabic')),
]
USE_I18N = True
USE_L10N = True
LOCALE_PATHS = [BASE_DIR / 'locale']

MIDDLEWARE = [
    'django.middleware.security.SecurityMiddleware',
    'django.contrib.sessions.middleware.SessionMiddleware',
    'django.middleware.locale.LocaleMiddleware',  # <-- ordre critique
    'django.middleware.common.CommonMiddleware',
    'django.middleware.csrf.CsrfViewMiddleware',
    'django.contrib.auth.middleware.AuthenticationMiddleware',
    'django.contrib.messages.middleware.MessageMiddleware',
]
\end{lstlisting}

Le template de base bascule la direction d'écriture en RTL pour l'arabe~:

\begin{lstlisting}[language=HTML,caption={templates/base.html (extrait)}]
<html lang="{{ LANGUAGE_CODE }}"
      dir="{% if LANGUAGE_BIDI %}rtl{% else %}ltr{% endif %}">
\end{lstlisting}

\subsection{Modèle d'un domaine clé --- la campagne}

\begin{lstlisting}[language=Python,caption={apps/pipeline/models.py (extrait)}]
from django.db import models
from django.utils.translation import gettext_lazy as _

class Campaign(models.Model):
    class Status(models.TextChoices):
        DRAFT = 'draft', _('Draft')
        RUNNING = 'running', _('Running')
        DONE = 'done', _('Done')

    product = models.ForeignKey(
        'pipeline.Product', on_delete=models.CASCADE,
        related_name='campaigns')
    name = models.CharField(_('Name'), max_length=120)
    start_date = models.DateField()
    end_date = models.DateField(null=True, blank=True)
    budget = models.DecimalField(max_digits=12, decimal_places=2)
    status = models.CharField(max_length=10, choices=Status.choices,
                              default=Status.DRAFT)

    def total_spent(self):
        return self.expenses.aggregate(s=models.Sum('amount'))['s'] or 0

    def total_revenue(self):
        return self.revenues.aggregate(s=models.Sum('amount'))['s'] or 0

    def roi(self):
        spent = self.total_spent()
        if not spent:
            return None
        return (self.total_revenue() - spent) / spent * 100
\end{lstlisting}

\subsection{Intégration de l'API DeepSeek}

Le service est volontairement minimaliste~: pas de SDK propriétaire, juste
un appel HTTP. La compatibilité avec le format OpenAI permet d'utiliser le
même schéma de messages.

\begin{lstlisting}[language=Python,caption={apps/chatbot/services.py (extrait)}]
import os, requests
from django.conf import settings
from .context import build_crm_context

DEEPSEEK_URL = settings.DEEPSEEK_BASE_URL + '/chat/completions'

def ask_deepseek(user, user_message, history, language='en'):
    """Send a message to DeepSeek with live CRM context injected."""
    facts = build_crm_context(user)  # top campaigns, ROI, best angle...

    system_prompt = (
        f"You are a marketing assistant. Reply in language={language}. "
        f"Use ONLY these facts about the user's CRM:\n{facts}"
    )
    payload = {
        'model': settings.DEEPSEEK_MODEL,
        'messages': [
            {'role': 'system', 'content': system_prompt},
            *history,
            {'role': 'user', 'content': user_message},
        ],
        'stream': False,
        'temperature': 0.4,
    }
    headers = {'Authorization': f'Bearer {settings.DEEPSEEK_API_KEY}'}
    r = requests.post(DEEPSEEK_URL, json=payload, headers=headers, timeout=30)
    r.raise_for_status()
    return r.json()['choices'][0]['message']['content']
\end{lstlisting}

\subsection{Modèle prédictif --- catégorisation du potentiel}

Le pipeline d'apprentissage est isolé dans \texttt{ml/train.py}, ce qui
permet de ré-entraîner sans toucher à l'application web.

\begin{lstlisting}[language=Python,caption={ml/train.py (extrait)}]
import joblib, pandas as pd
from sklearn.ensemble import RandomForestClassifier
from sklearn.model_selection import train_test_split
from sklearn.metrics import classification_report

df = pd.read_csv('ml/seed_data.csv')
X = pd.get_dummies(df.drop(columns=['potential']))
y = df['potential']  # Low / Medium / High

Xtr, Xte, ytr, yte = train_test_split(X, y, test_size=0.2, random_state=42)
clf = RandomForestClassifier(n_estimators=200, random_state=42)
clf.fit(Xtr, ytr)
print(classification_report(yte, clf.predict(Xte)))

joblib.dump({'model': clf, 'columns': X.columns.tolist()}, 'ml/model.pkl')
\end{lstlisting}

Le service d'inférence charge le modèle une seule fois au démarrage~:

\begin{lstlisting}[language=Python,caption={apps/prediction/services.py (extrait)}]
import joblib, pandas as pd
from functools import lru_cache

@lru_cache(maxsize=1)
def _bundle():
    return joblib.load('ml/model.pkl')

def predict_potential(features: dict):
    b = _bundle()
    X = pd.DataFrame([features]).reindex(columns=b['columns'], fill_value=0)
    proba = b['model'].predict_proba(X)[0]
    classes = b['model'].classes_
    idx = proba.argmax()
    return classes[idx], float(proba[idx])
\end{lstlisting}

\subsection{Pipeline Kanban --- HTMX et glisser-déposer}

Le drag-and-drop ne nécessite ni librairie tierce ni framework JS~: la
balise HTMX \texttt{hx-post} renvoie le nouveau statut au backend, qui
répond avec le fragment HTML mis à jour.

\begin{lstlisting}[language=HTML,caption={templates/pipeline/kanban.html (extrait)}]
<div class="kanban-col" data-status="active"
     hx-post="{% url 'pipeline:move' %}"
     hx-trigger="drop"
     hx-vals='{"status": "active"}'
     hx-swap="outerHTML">
  {% for product in active_products %}
    <article draggable="true" data-id="{{ product.id }}">
      {{ product.name }}
    </article>
  {% endfor %}
</div>
\end{lstlisting}
